%% ma: L10-B02-0025
%% lop: 10
%% chuong: 1
%% bai: 2
%% dang: 10.2.2
%% loai: TN4
%% muc: TH
%% dap_an: "D"
%% nguon: "Ảnh giáo viên gửi 10/10/2026 (ThS. Nguyễn Hoàng Việt, Chương 2 Tập hợp), Đề 2 (đề thứ hai, không ghi tiêu đề), Phần 1 Câu 1"
%% kien_thuc: "Bài 2 (tập hợp và các phép toán trên tập hợp)"
%% trang_thai: cho-duyet
%% ly_do: "Dạng toán mới đề xuất 10.2.2: xác định tập hợp bằng cách liệt kê phần tử (giải phương trình, bất phương trình trong N, Z, Q, R)"
\begin{ex}
Cho tập hợp $A=\{x\in\mathbb N\mid(x^3-9x)(2x^2-5x+2)=0\}$. Tập $A$ được viết theo kiểu liệt kê là
\choice{$\{2;3\}$}{$\left\{-3;0;\dfrac12;2;3\right\}$}{$\{-3;0;2;3\}$}{\True $\{0;2;3\}$}
\loigiai{$x^3-9x=x(x-3)(x+3)=0\Leftrightarrow x\in\{0;3;-3\}$; $2x^2-5x+2=(2x-1)(x-2)=0\Leftrightarrow x\in\left\{\dfrac12;2\right\}$. Chỉ các nghiệm tự nhiên được nhận: $A=\{0;2;3\}$.}
\end{ex}
